\documentclass[11pt]{article}
\usepackage[a4paper,margin=18mm]{geometry}
\usepackage{fontspec}
\setmainfont{Latin Modern Roman}
\usepackage{amsmath,amssymb,amsthm,amscd,graphicx,color}
\usepackage{xurl}
\usepackage[unicode,hidelinks]{hyperref}
\allowdisplaybreaks
\setlength\emergencystretch{3em}
\numberwithin{equation}{section}

\newtheorem{Theorem}{Theorem}[section]
\newtheorem{Lemma}[Theorem]{Lemma}
{ \theoremstyle{definition}
\newtheorem{Definition}[Theorem]{Definition}
\newtheorem{Example}[Theorem]{Example}}
\let\originalcedilla\c
\renewcommand{\c}[1]{{\fontspec{DejaVu Serif}\originalcedilla{#1}}}

\makeatletter\newlabel{5}{{2.1}{}{Original source locator}{}{}}\makeatother
\makeatletter\newlabel{6}{{2.2}{}{Original source locator}{}{}}\makeatother
\begin{document}
\begin{center}\Large Principality of algebraic equivariant fusion for Hopf algebras with bijective antipode, principal comodule algebras and two jointly surjective C-star quotient maps\end{center}
\noindent\textbf{Stable ID:} 2013\par
\noindent\textbf{Authors:} Ludwik Dabrowski; Tom Hadfield; Piotr M. Hajac\par
\noindent\textbf{Original work:} Equivariant Join and Fusion of Noncommutative Algebras\par
\noindent\textbf{Version:} Official SIGMA 11 (2015) article 082, DOI 10.3842/SIGMA.2015.082; journal PDF and native TeX acquired 2026-10-01, exact bytes pinned by SHA256\par
\noindent\textbf{Selected task scope:} Full Theorem5.3 in the algebraic Section5 setting: equivariant fusion P fusion\_delta,pi H is principal under the standing bijective-antipode and principal-right-H-comodule assumptions, with unital C-star algebra C and jointly surjective quotient map C->C/J1 directsum C/J2. No claim that arbitrary non-surjective maps, generic diagonal coactions or a C-star completion satisfy this theorem.\par
\noindent\textbf{Difficulty:} H2: The construction combines the original strong connection with an antipode-controlled trivial connection through two square-root cutoff elements obtained by C-star functional calculus. The proof verifies both quotient boundary constraints and the canonical splitting, then derives Hopf-Galois bijectivity and equivariant projectivity from a nonunital connection criterion; it covers fusion without a pullback join decomposition.\par
\noindent\textbf{Editorial boundary note (not author text):} Full original Theorem 5.3 and complete new square-root strong-connection construction are supplied with the nonunital connection criterion and full fusion coaction restriction and boundary-corestriction proof. H has bijective antipode, P is principal in both Hopf-Galois and equivariant-projectivity senses, C is a unital complex C-star algebra and the direct sum of the two canonical ideal quotient maps is surjective. Tensor products in the selected fusion are algebraic. Every boundary calculation and the final splitting identity is retained; bicolinearity is the source author argument from the input connection. Prior strong-connection equivalence and functional calculus remain named standard background. No ordinary-join example, analytic completion or conjectural nontriviality claim is selected.\par
\noindent\textbf{Source:} \url{https://sigma-journal.com/2015/082/}\quad DOI: \url{https://doi.org/10.3842/SIGMA.2015.082}\par
\noindent\textbf{License and attribution:} Copyright retained by the named authors. Published by SIGMA under CC BY 4.0: \url{https://creativecommons.org/licenses/by/4.0/}. Source: \url{https://sigma-journal.com/about.html}. Adaptation consists of standalone excerpt formatting, cover and numbering restoration; original author language and mathematical text are retained.\par
\noindent\textbf{Typesetting adaptation (not author text):} The native cedilla accent in cited author names uses the portable DejaVu Serif font, because Latin Modern Roman lacks its combining cedilla glyph. Accent commands and bibliography text remain unchanged.\par
\noindent\textbf{External original locator 5:} Equation (2.1) in the original article is a motivational classical join equivalence relation. The selected algebraic fusion and its proof are explicitly defined below.\par
\noindent\textbf{External original locator 6:} Equation (2.2) in the original article is the motivational classical gauging map; no unbundled new proof obligation is assigned to it.\par
\medskip\hrule\medskip\noindent\textbf{Author text begins below.}\par
\setcounter{section}{2}
\setcounter{subsection}{0}
\setcounter{subsubsection}{0}
\setcounter{Theorem}{0}
\setcounter{equation}{2}
% Original source lines 173--321
\section{Strong connection and principal comodule algebra}\label{sc}

Let $H$ be a Hopf algebra with coproduct~$\Delta$, counit~$\varepsilon$ and bijective
antipode~$S$. Next, let \mbox{$\delta\colon P\to P\otimes H$} be a coaction making $P$ a right
\mbox{$H$-comodule} algebra.
We shall frequently use
the Heyneman--Sweedler notation (with the summation sign suppressed) for coproduct and coaction:
\begin{gather*}
\Delta (h)=:h_{(1)}\otimes h_{(2)} ,\qquad
\delta(p)=:p_{(0)}\otimes p_{(1)} .
\end{gather*}
\begin{Definition}[\cite{hkmz11}]
Let $P$ be a right comodule algebra over a Hopf algebra
$H$ with bijective antipode, and let
\begin{gather*}
B:=P^{\operatorname{co}H} := \{ p \in P \, | \, \delta (p) = p \otimes 1 \}.
\end{gather*}
be the coaction-invariant
subalgebra. The comodule algebra $P$
is called {\em principal} if the following conditions are satisf\/ied:
\begin{enumerate}\itemsep=0pt
\item[1)]
the coaction of $H$ is Hopf--Galois, that is, the map
\begin{gather*}
\operatorname{can}_P \colon \  P \underset{B}{\otimes} P \longrightarrow P \otimes H,\qquad
	p \otimes q \longmapsto pq_{(0)} \otimes q_{(1)} ,
\end{gather*}
(called the canonical map) is bijective,
\item[2)]
the comodule algebra $P$ is right
$H$-equivariantly projective as a left  $B$-module, i.e.,
 there exists
 a right $H$-colinear and left $B$-linear
splitting of the multiplication map
$B \otimes P \rightarrow P$.
\end{enumerate}
\end{Definition}

\begin{Definition}[\cite{bh04}]\label{strong}
Let $ H$ be a Hopf algebra with  bijective antipode.
A \emph{strong connection $\ell$} on a right $H$-comodule algebra $P$ is a unital linear map $\ell\colon  H
\rightarrow  P \otimes  P$ satisfying
the following bicolinearity (i.e., left and right colinearity) and splitting conditions:
\begin{enumerate}\itemsep=0pt
\item[1)]
$(\operatorname{id}\otimes \delta) \circ
\ell = (\ell \otimes \operatorname{id}) \circ \Delta$,
$(\delta^L \otimes \operatorname{id}) \circ
\ell = (\operatorname{id} \otimes \ell) \circ
\Delta$, where
$\delta^L :=(S^{-1}\otimes\operatorname{id})\circ\mathrm{f\/lip}\circ\delta$;
\item[2)]
$\widetilde{\operatorname{can}_P} \circ \ell=1 \otimes \operatorname{id}$, where
$\widetilde{\operatorname{can}_P}\colon  P\otimes  P\ni p\otimes q\mapsto (p\otimes 1)\delta(q)\in  P\otimes H$.
\end{enumerate}
\end{Definition}

We will use the  Heyneman--Sweedler-type notation
$
\ell(h)=:\ell(h)^{\langle 1 \rangle}\otimes
 \ell(h)^{\langle 2 \rangle}
$
 with the summation sign suppressed.
 One can easily prove (see \cite[p.~599]{hkmz11} and references therein) that
a~comodule algebra is principal if and only if it admits a strong
connection.
Here we need the following slight generalization of this fact:
\begin{Lemma}\label{general}
Let $H$ be a Hopf algebra with bijective antipode. Then a right $H$-comodule algebra~$P$ is principal if and
only if it admits a~$($not necessarily unital$)$ linear map $\ell \colon H \rightarrow  P \otimes  P$ satisfying the
bicolinearity and splitting conditions of Definition~{\rm \ref{strong}}.
\end{Lemma}

\begin{proof}
Assume that $\ell \colon H \rightarrow  P \otimes  P$ is a linear map satisfying the bicolinearity and splitting
conditions.
Let $\pi_B\colon P\otimes P\to P\otimes_BP$ be the canonical surjection. Def\/ine
\begin{gather*}
L\colon \ P\otimes H\ni p\otimes h\longmapsto \pi_B\big(p \ell(h)^{\langle 1\rangle}\otimes\ell(h)^{\langle 2\rangle}\big)\in P\underset{B}{\otimes}P.
\end{gather*}
It follows immediately from the  splitting property of $\ell$ that
\begin{gather*}
\operatorname{can}_P(L(p\otimes h))= p \, \widetilde{\operatorname{can}_P}(\ell(h))=p\otimes h.
\end{gather*}
Applying $\operatorname{id}\otimes\varepsilon$ to the splitting condition for $\ell$, we obtain $m\circ\ell=\varepsilon$, where $m$ is
the multiplication map on~$P$.
Combining it with the left colinearity of $\ell$, which implies that
\begin{gather*}
\delta\big(q_{(0)}\ell(q_{(1)})^{\langle 1\rangle}\big)\otimes\ell(q_{(1)})^{\langle 2\rangle}
=q_{(0)}\ell(q_{(3)})^{\langle 1\rangle}\otimes q_{(1)}S(q_{(2)})\otimes\ell(q_{(3)})^{\langle 2\rangle}\\
\hphantom{\delta\big(q_{(0)}\ell(q_{(1)})^{\langle 1\rangle}\big)\otimes\ell(q_{(1)})^{\langle 2\rangle}}{}
=q_{(0)}\ell(q_{(1)})^{\langle 1\rangle}\otimes 1\otimes\ell(q_{(1)})^{\langle 2\rangle},
\end{gather*}
we obtain
\begin{gather*}
L(\operatorname{can}_P(p\otimes q)) = \pi_B \big(pq_{(0)} \ell(q_{(1)})^{\langle 1\rangle}
 \otimes  \ell(q_{(1)})^{\langle 2\rangle}\big)
 = \pi_B \big(p \otimes q_{(0)} \ell(q_{(1)})^{\langle 1\rangle} \ell(q_{(1)})^{\langle 2\rangle}\big)
 \\
\hphantom{L(\operatorname{can}_P(p\otimes q))}{}
=  \pi_B(p\otimes q).
\end{gather*}
Finally, the bicolinearity of $\ell$ implies that the formula $s(p):= p_{(0)}\ell(p_{(1)})$ def\/ines a left $B$-linear
right $H$-colinear splitting
of the multiplication map $B\otimes P\to P$. The reverse implication (principality $\Rightarrow$ existence of a~bicolinear~$\ell$ with the splitting property)
follows from~\cite[Lemma~2.2]{bh04}.
\end{proof}


\section[Join and fusion of unital $C^*$-algebras]{Join and fusion of unital $\boldsymbol{C^*}$-algebras}


Let $\otimes_{\min}$ denote the spatial (minimal) tensor product.
Due to the fact that minimal tensor products preserve injections
 (e.g., see \cite[Proposition~4.22]{t-m79} or
\cite[Section~1.3]{w-s94}),
for any unital $C^*$-algebras $A_1$ and $A_2$ the natural maps
\begin{gather*}
A_1\ni a\longmapsto a\otimes 1\in A_1 \underset{\min}{\otimes} A_2\qquad \text{and}
\qquad  A_2\ni a\longmapsto 1\otimes a\in A_1 \underset{\min}{\otimes} A_2
\end{gather*}
are injective. This allows us to view $A_1$ and $A_2$ as subalgebras of
$A_1 \otimes_{\min} A_2$. We use this natural identif\/ication in what follows.
\begin{Definition}\label{gjoin}
Let $A_1$ and $A_2$  be unital $C^*$-algebras, and
let $C$ be a unital $C^*$-algebra equipped with two $C^*$-ideals $J_1$ and $ J_2$
such that the direct sum
$
\pi_1\oplus\pi_2\colon C\rightarrow (C/J_1) \oplus (C/J_2)
$
of the canonical surjections $\pi_i\colon C\to C/J_i$, $i\in \{1,2\}$, is surjective.
We call the unital $C^*$-algebra
\begin{gather*}
A_1 \underset{\pi}{\circledast} A_2 :=
\big\{x\in C \underset{\min}{\otimes} A_1 \underset{\min}{\otimes} A_2 \,
\big|\,
(\pi_i\otimes \operatorname{id})(x) \in
(C/J_i)\otimes_{\min}A_i
, \, i\in\{1,2\}
\big\}
\end{gather*}
the \emph{fusion} $C^*$-algebra of $A_1$ and~$A_2$ over~$C$.
When $C$ is the $C^*$-algebra $C([0,1])$ of all continuous complex-valued
functions on the unit interval and
$\pi_1:=\mathrm{ev}_1$ and $\pi_2:=\mathrm{ev}_0$ are
the evaluation maps respectively at $1$ and $0$,
we call $A_1 \circledast_\pi A_2$ the \emph{join} $C^*$-algebra of $A_1$ and~$A_2\,$, and denote it by
 $A_1 \circledast A_2$.
\end{Definition}
\setcounter{section}{4}
\setcounter{subsection}{0}
\setcounter{subsubsection}{0}
\setcounter{Theorem}{1}
\setcounter{equation}{0}
% Original source lines 340--473
\section[Join and fusion of a principal $H$-comodule algebra with $H$]{Join and fusion of a principal $\boldsymbol{H}$-comodule algebra with $\boldsymbol{H}$}

Since a diagonal coaction is not in general an algebra homomorphism,
to obtain an equivariant version of our noncommutative fusion and join constructions,
we need to modify  Def\/inition~\ref{gjoin}
in the spirit of~\eqref{5},~\eqref{6}.
To avoid analytical complications that
are tackled in~\cite{ddhw}, we also need to change the setting from $C^*$-algebraic to algebraic. In particular,
we replace $\otimes_{\min}$ by~$\otimes$.

\begin{Definition}\label{pjoin}
Let $\delta\colon P\to P\otimes H$ be a right coaction making $P$ a comodule algebra, and
let~$C$ be a unital $C^*$-algebra equipped with two $C^*$-ideals $J_1$ and~$ J_2$
such that the direct sum
\begin{gather*}
\pi_1\oplus\pi_2\colon \ C\longrightarrow (C/J_1) \oplus (C/J_2)
\end{gather*}
of the canonical surjections $\pi_i\colon C\to C/J_i$, $i\in \{1,2\}$, is surjective. We call the algebra
\begin{gather*}
P \underset{\pi}{\overset{\delta}{\circledast}} H :=
\big\{x\in C \otimes P \otimes H \,
\big|\, \\
\hphantom{P \underset{\pi}{\overset{\delta}{\circledast}} H :=\big\{}{}
(\pi_1\otimes \operatorname{id})(x) \in (C/J_1)\otimes \delta (P) ,\,
(\pi_2\otimes \operatorname{id})(x) \in (C/J_2) \otimes \mathbb{C} \otimes H\big\}
\end{gather*}
the \emph{equivariant fusion} algebra of $P$ and~$H$ over~$C$.
When $C:=C([0,1])$  and
$\pi_1:=\mathrm{ev}_1$, $\pi_2:=\mathrm{ev}_0$,
we call $P \circledast^\delta_\pi H$ the \emph{equivariant join} algebra of $P$ and~$H$,
and denote it by $P \circledast^\delta H$.
\end{Definition}

\begin{Lemma}
The coaction $\operatorname{id}\otimes\operatorname{id}\otimes\Delta\colon
C\otimes P\otimes H\to C\otimes P\otimes H\otimes H$
restricts and corestricts to $\Delta_{P\circledast{}^\delta_\pi H}\colon P{\circledast}^\delta_\pi   H
\to (P{\circledast}^\delta_\pi H)\otimes H$ making $P{\circledast}^\delta_\pi H$
a right $H$-comodule algebra.
\end{Lemma}

\begin{proof}
Note f\/irst that for $i\in\{1,2\}$, we have
\begin{gather*}
\big((\pi_i\otimes\operatorname{id}\otimes\operatorname{id})\circ(\operatorname{id}\otimes\operatorname{id}\otimes\Delta)\big)\!\bigg(\sum_j f_j\otimes p_j\otimes h_j\bigg)
=(\operatorname{id}\otimes\operatorname{id}\otimes\Delta)\!\bigg(\sum_j \pi_i(f_j) \otimes p_j\otimes h_j\bigg).
\end{gather*}
Assume now that $\sum_j f_j\otimes p_j\otimes h_j\in P\circledast^\delta_\pi H$.
 Then the above tensor belongs to
\begin{gather*}
(C/J_1)\otimes \big((\operatorname{id}\otimes\Delta)(\delta(P))\big)=
(C/J_1)\otimes \big((\delta\otimes\operatorname{id})(\delta(P))\big)
\subseteq(C/J_1)\otimes \delta(P)\otimes H
\end{gather*}
for $i=1$ and to $(C/J_2)\otimes \mathbb{C}\otimes H\otimes H$ for~$i=2$.
\end{proof}

The main result of this paper is:
\begin{Theorem}\label{main}
Let $C$ be a unital $C^*$-algebra equipped with two $C^*$-ideals $J_1$ and $ J_2$
such that the direct sum
$
\pi_1\oplus\pi_2\colon C\rightarrow (C/J_1) \oplus (C/J_2)
$
of the canonical surjections $\pi_i\colon C\to C/J_i$, $i\in \{1,2\}$, is surjective.
Then, for any principal right $H$-comodule algebra~$P$,
 the equivariant fusion right $H$-comodule algebra $P{\circledast}^\delta_\pi H$ is principal.
\end{Theorem}

\begin{proof}
Note f\/irst that, due to  the surjectivity of $\pi_1\oplus\pi_2$, we can choose $x\in C$  such that
\begin{gather*}
(\pi_1\oplus\pi_2)(x)=(1,0).
\end{gather*}
Since $\{0, 1\}\subseteq \operatorname{spec}(x^*x)$, we can
 take a continuous function $f\colon \operatorname{spec}(x^*x)\to [0,1]$
such that $f(0)=0$ and $f(1)=1$, and def\/ine $t:=f(x^*x)\in C$.
Then there exist elements $\sqrt{t},\sqrt{1-t}\in C$ enjoying the properties
$\pi_2(\sqrt{t})=0=\pi_1(\sqrt{1-t})$.

Next, let
$\ell \colon H{\to} P\otimes P$  be a strong connection on $P$.
Then the bicolinearity of~$\ell$ implies that the linear map
\begin{gather*}
\tilde\ell \colon \  H \longrightarrow \big( C \otimes P \otimes H\big) \otimes \big( C \otimes P \otimes H\big),\\
\tilde\ell(h):= {\sqrt t} \otimes  \ell (h_{(2)})^{\langle 1\rangle}\otimes S(h_{(1)})
\otimes {\sqrt t} \otimes  \ell (h_{(2)})^{\langle 2\rangle}  \otimes h_{(3)}
\\
\hphantom{\tilde\ell(h):=}{}+ \sqrt{1-t} \otimes 1 \otimes S(h_{(1)}) \otimes \sqrt{1-t}
\otimes 1\otimes h_{(2)},
\end{gather*}
corestricts to $(P\circledast^\delta_\pi H)\otimes  (P\circledast^\delta_\pi H)$.
Indeed,
\begin{gather*}
 (\pi_1\otimes\operatorname{id}\otimes\operatorname{id})\otimes(\operatorname{id}\otimes\operatorname{id}\otimes\operatorname{id})(\widetilde{\ell}(h))\\
 \qquad{} =
\big([1]\otimes\delta\big(\ell(h_{(1)})^{\langle 1\rangle}\big)\big)\otimes
\big(\sqrt{t}\otimes \ell(h_{(1)})^{\langle 2\rangle}\otimes h_{(2)}\big)
\in \big(C/J_1\otimes\delta(P)\big)\otimes\big(C\otimes P\otimes H\big),\\
 (\pi_2\otimes\operatorname{id}\otimes\operatorname{id})\otimes(\operatorname{id}\otimes\operatorname{id}\otimes\operatorname{id})(\widetilde{\ell}(h))\\
 \qquad{} =
\big([1]\otimes 1\otimes S(h_{(1)}\big)\otimes
\big(\sqrt{1-t}\otimes 1\otimes h_{(2)}\big)
\in \big(C/J_2\otimes\mathbb{C}\otimes H\big)\otimes\big(C\otimes P\otimes H\big).
\end{gather*}
Much in the same way,
\begin{gather*}
 (\operatorname{id}\otimes\operatorname{id}\otimes\operatorname{id})\otimes(\pi_1\otimes\operatorname{id}\otimes\operatorname{id})(\widetilde{\ell}(h))
\in \big(C\otimes P\otimes H\big)\otimes\big(C/J_1\otimes\delta(P)\big),\\
 (\operatorname{id}\otimes\operatorname{id}\otimes\operatorname{id})\otimes(\pi_2\otimes\operatorname{id}\otimes\operatorname{id})(\widetilde{\ell}(h))
\in\big(C\otimes P\otimes H\big) \otimes\big(C/J_2\otimes\mathbb{C}\otimes H\big).
\end{gather*}

The bicolinearity of the thus corestricted $\tilde\ell$ is evident. Finally, taking
advantage of the right colinearity and the splitting property of~$\ell$, we check that
\begin{gather*}
 \widetilde{\operatorname{can}_{P\circledast^\delta_\pi H}}(\tilde\ell (h)) =
 t \otimes  \ell (h_{(2)})^{\langle 1\rangle}\ell (h_{(2)})^{\langle 2\rangle}
\otimes S(h_{(1)})  h_{(3)}\otimes h_{(4)}\\
\hphantom{\widetilde{\operatorname{can}_{P\circledast^\delta_\pi H}}(\tilde\ell (h)) =}{}
+ (1-t) \otimes 1 \otimes S(h_{(1)})h_{(2)} \otimes h_{(3)}
 \\
\hphantom{\widetilde{\operatorname{can}_{P\circledast^\delta_\pi H}}(\tilde\ell (h))}{} =  t \otimes  \ell (h_{(2)})^{\langle 1\rangle}\ell (h_{(2)})^{\langle 2\rangle}\!{}_{(0)} \otimes S(h_{(1)})
\ell (h_{(2)})^{\langle 2\rangle}\!{}_{(1)}\otimes h_{(3)}
+ (1-t) \otimes 1 \otimes 1 \otimes h
 \\
\hphantom{\widetilde{\operatorname{can}_{P\circledast^\delta_\pi H}}(\tilde\ell (h))}{} =  t \otimes  1 \otimes S(h_{(1)})  h_{(2)}\otimes h_{(3)}
+ (1-t) \otimes 1 \otimes 1 \otimes h
 \\
\hphantom{\widetilde{\operatorname{can}_{P\circledast^\delta_\pi H}}(\tilde\ell (h))}{}
 =  1 \otimes  1 \otimes 1 \otimes h.
\end{gather*}
Now the claim follows from Lemma~\ref{general}.
\end{proof}
\begin{thebibliography}{99}
\bibitem[3]{bh04}
Brzezi{\'n}ski T., Hajac P.M., The {C}hern--{G}alois character, \href{http://dx.doi.org/10.1016/j.crma.2003.11.009}{\textit{C.~R.
  Math. Acad. Sci. Paris}} \textbf{338} (2004), 113--116,
  \href{http://arxiv.org/abs/math.KT/0306436}{math.KT/0306436}.

\bibitem[5]{ddhw}
D\c{a}browski L., De~Commer K., Hajac P.M., Wagner E., Noncommutative bordism
  of free actions of compact quantum groups on unital $C^*$-algebras, {i}n
  preparation.

\bibitem[7]{hkmz11}
Hajac P.M., Kr{\"a}hmer U., Matthes R., Zieli{\'n}ski B., Piecewise principal
  comodule algebras, \href{http://dx.doi.org/10.4171/JNCG/88}{\textit{J.~Noncommut. Geom.}} \textbf{5} (2011), 591--614,
  \href{http://arxiv.org/abs/0707.1344}{arXiv:0707.1344}.

\bibitem[11]{t-m79}
Takesaki M., Theory of operator algebras.~{I}, \href{http://dx.doi.org/10.1007/978-1-4612-6188-9}{Springer-Verlag}, New York~--
  Heidelberg, 1979.

\bibitem[12]{w-s94}
Wassermann S., Exact {$C^*$}-algebras and related topics, \textit{Lecture Notes
  Series}, Vol.~19, Seoul National University, Research Institute of
  Mathematics, Global Analysis Research Center, Seoul, 1994.

\end{thebibliography}
\end{document}
